\par\Needspace{9\baselineskip}\bigskip
\pdfbookmark[3]{Exposition, [22]}{pallas-II-article22}
\begin{center}[22.]\end{center}
\begingroup
\setlength{\lineskip}{4pt}\setlength{\lineskiplimit}{4pt}
\noindent\hspace*{\parindent}The instantaneous variations of each element are evidently products of \(\mathrm{d}t\) and functions of time and of the elements of the \SourceLeaf{472}perturbed planet and of the perturbing planet. In the first approximation, all these elements are supposed constant, and time enters only through the values of the mean longitudes of the two planets, which are
\[
\begin{aligned}
L&=N+nt\\
L'&=N'+n't,
\end{aligned}
\]
where we distinguish by primes all quantities relating to the perturbing planet. In this way the instantaneous variation of each element takes the form \(\mathrm{d}t.(\text{function of }L\text{ and }L')\). It is also evident that this function will be \textit{periodic} in \(L\text{ and }L'\), because it will [not] change in value if \(L\text{ and }L'\) are varied by \(360^{0}\) or a multiple of \(360^{0}\). Functions of this kind are known to be expressible as a series of terms of the form \(k\cos(iL+i'L')\) or \(k\sin(iL+i'L')\), where \(k\) is a constant coefficient, \(i\) runs through all positive integers (including zero), and \(i'\) through all positive or negative integers. On integration, the term for which \(i=i'=0\) will give \(kt\), the secular variation of the element. The other terms will give periodic parts of the form \(\dfrac{k}{in+i'n'}\sin(iL+i'L')\). If any terms for which \(in+i'n'=0\) were found, they would have to be excepted and would add \(kt\cos(iL+i'L')\) to the secular variation. It is evident, however, that such terms can occur only when the mean motions are in a rational ratio. Finally, a constant must be added.
\par\endgroup
